% ============================================================
%  main.tex  —  complete standalone file for Overleaf
%  Compiler: XeLaTeX  (Menu → Compiler → XeLaTeX)
%  Upload alongside:  ti_projection_curve.png
% ============================================================
\documentclass[11pt,a4paper]{article}

% ── Unicode & font (XeLaTeX + Overleaf handle Ethiopic natively) ──────────────
\usepackage{fontspec}
\usepackage{polyglossia}
\setmainlanguage{english}
\newfontfamily\ethiopicfont[Script=Ethiopic]{Noto Sans Ethiopic}

% ── Mathematics ───────────────────────────────────────────────────────────────
\usepackage{amsmath}
\usepackage{amssymb}

% ── Tables ────────────────────────────────────────────────────────────────────
\usepackage{booktabs}
\usepackage{array}
\usepackage{multirow}

% ── Figures ───────────────────────────────────────────────────────────────────
\usepackage{graphicx}

% ── Layout ────────────────────────────────────────────────────────────────────
\usepackage[margin=2.5cm]{geometry}
\usepackage{microtype}
\usepackage{parskip}
\usepackage{setspace}
\onehalfspacing

% ── Hyperlinks ────────────────────────────────────────────────────────────────
\usepackage[colorlinks=true, linkcolor=blue, citecolor=blue, urlcolor=blue]{hyperref}

% ── Tigrigna inline helper ────────────────────────────────────────────────────
% XeLaTeX renders Ethiopic Unicode directly; \ti{} just switches font explicitly
\newcommand{\ti}[1]{{\ethiopicfont #1}}


% ══════════════════════════════════════════════════════════════════════════════
\begin{document}
% ══════════════════════════════════════════════════════════════════════════════

\title{\textbf{Layer-wise Gender Bias in Multilingual Transformers:\\
       Evidence from Tigrigna, a Low-Resource Language}}

\author{[Author Names] \\ [Institution] \\ \texttt{[email]}}
\date{}
\maketitle

% ── Abstract ──────────────────────────────────────────────────────────────────
\begin{abstract}
We investigate gender bias in multilingual transformer models as reflected in
the contextual representations of Tigrigna profession words, probing all
transformer layers of five models — XLM-RoBERTa-base (13 layers),
XLM-RoBERTa-large (25 layers), XLM-V-base (13 layers),
mDeBERTa-v3-base (13 layers), and mBERT (13 layers) — using two independent
bias estimators. We find that mBERT maps all Tigrigna tokens to \texttt{[UNK]},
producing degenerate identical representations and zero measurable bias signal —
establishing vocabulary coverage as a prerequisite for multilingual bias
measurement. XLM-RoBERTa-base, XLM-RoBERTa-large, and XLM-V-base all produce
near-static representations and consistent \emph{male}-leaning bias for 25--26
of 30 professions, with layer-wise peak projections of $+2.99$, $+3.46$, and
$+2.63$ respectively. The most striking finding comes from mDeBERTa-v3-base:
its Disentangled Attention architecture produces a complete reversal, with
28 of 30 professions \emph{female}-leaning — the opposite direction from all
XLM-R family models. XLM-RoBERTa-large further exhibits a two-phase double-peak
bias structure (layers~13 and~23). Comparing static and contextual embeddings
reveals two qualitatively distinct mechanisms: \emph{embedding-level bias} and
\emph{attention-level bias}. ProjBias and CosDiff estimators agree near-perfectly
across all layers in all four meaningful models (Spearman $\rho \geq 0.970$,
$p < 10^{-17}$), and batch-stability validation confirms estimates are fully
converged and sampling-independent.
\end{abstract}

% ══════════════════════════════════════════════════════════════════════════════
\section{Methodology}
% ══════════════════════════════════════════════════════════════════════════════

\subsection{Language and Corpus Construction}

We focus on Tigrigna (\texttt{ti}), a Semitic language spoken primarily in
Eritrea and the Tigray region of Ethiopia, written in Ge'ez (Ethiopic) script.
Tigrigna is a genuinely low-resource language: no substantial pretrained
monolingual model exists, and its representation in multilingual corpora is
sparse.

Because no sufficiently large, profession-annotated Tigrigna corpus is publicly
available, we constructed a synthetic corpus of \textbf{14{,}960 sentences}.
The corpus was built using controlled sentence templates, each instantiating one
of 30 target profession words in a short declarative context. Templates were
designed to co-occur the profession word with gender-marked anchor terms
(pronouns, kinship terms, and generic gender nouns) drawn from Tigrigna
morphology. Every profession word appears in exactly 499 sentences, distributed
across male-gendered, female-gendered, and gender-neutral contexts. Each
sentence was automatically labelled \texttt{M}, \texttt{F}, or \texttt{N} by
checking for the presence of known gender marker tokens
(see Table~\ref{tab:anchors}).

\begin{table}[h]
\centering
\small
\begin{tabular}{llll}
\toprule
\textbf{Gender} & \textbf{Type} & \textbf{Tigrigna} & \textbf{English gloss} \\
\midrule
Male   & Pronoun & ንሱ, ንዕኡ, ናቱ             & he, him, his \\
Male   & Kinship & ኣቦ, ሓው, ወዲ ሓው, በዓል ገዛ  & father, brother, nephew, husband \\
Male   & Generic & ሰብኣይ, ወዲ, ተባዕታይ, ኣቦሓጎ & man, boy/son, male, grandfather \\
\midrule
Female & Pronoun & ንሳ, ንዓኣ, ናታ            & she, her, her (poss.) \\
Female & Kinship & ኣደ, ሓፍቲ, ጓል ሓፍቲ, ሰበይቲ & mother, sister, niece, wife \\
Female & Generic & ጓል, ኣንስተይቲ, ዓባየይ       & girl/daughter, female, grandmother \\
\bottomrule
\end{tabular}
\caption{Gender anchor words used to build the gender subspace and classify
corpus sentences. 11 male and 10 female anchors are used.}
\label{tab:anchors}
\end{table}

\subsection{Models}

We evaluate five multilingual transformer models, all used strictly in inference
mode with frozen weights. All experiments run on CPU with full determinism.

\paragraph{XLM-RoBERTa-base}~\cite{conneau2020xlm} (\texttt{xlm-roberta-base}).
A 12-layer transformer with 768-dimensional hidden states, trained on 100
languages via masked language modelling on CommonCrawl. We extract hidden states
from all 13 layers (embedding layer~0 plus 12 transformer blocks).
Tokenisation uses a SentencePiece vocabulary of 250{,}000 pieces with
character-level fallback, giving broad Ethiopic script coverage.

\paragraph{XLM-RoBERTa-large} (\texttt{xlm-roberta-large}).
The large variant of the same family: 24 transformer layers, 1{,}024-dimensional
hidden states, 16 attention heads. Shares the identical SentencePiece vocabulary
and tokenisation scheme as the base model, so Tigrigna subword decomposition
is unchanged. We extract hidden states from all 25 layers.

\paragraph{XLM-V-base}~\cite{liang2023xlmv} (\texttt{facebook/xlm-v-base}).
An extended RoBERTa-style architecture with a dramatically enlarged SentencePiece
vocabulary of 1{,}000{,}000 pieces — four times that of XLM-RoBERTa — designed
to overcome vocabulary bottlenecks for low-resource languages. Architecture is
otherwise identical to XLM-RoBERTa-base: 12 transformer layers, 768-dimensional
hidden states. We extract 13 hidden states. Ethiopic script coverage is broad
due to the million-piece vocabulary; however, 2 of 30 professions
({\ethiopicfont ጠበቓ}, lawyer; {\ethiopicfont ነዳቓይ}, miner) tokenise differently
under this vocabulary from the substrings present in our corpus, causing zero
context matches and excluding them from analysis (28 professions scored).

\paragraph{mDeBERTa-v3-base}~\cite{he2021deberta} (\texttt{microsoft/mdeberta-v3-base}).
A 12-layer, 768-dimensional multilingual transformer using Disentangled
Attention, which encodes token content and relative positional information as
separate key/value pairs rather than combining them additively as in standard
self-attention. Trained on multilingual corpora with a SentencePiece tokeniser.
All 30 Tigrigna professions receive non-trivial subword decompositions, and we
extract 13 hidden states. As shown in Section~\ref{sec:deberta_reversal}, the
Disentangled Attention mechanism produces a fundamentally different gender bias
geometry compared to the XLM-R family.

\paragraph{mBERT} (\texttt{bert-base-multilingual-cased}).
A 12-layer, 768-dimensional BERT model trained on 104 languages using WordPiece
tokenisation with a ~120{,}000 token vocabulary. Included as a baseline to test
whether any multilingual model can measure Tigrigna bias.

For any word $w$ that tokenises into $k$ subword pieces $t_1, \ldots, t_k$,
we obtain its representation at layer $l$ as the mean of the $k$ hidden-state
vectors at the corresponding token positions:
\begin{equation}
  \mathbf{s}_l(w) = \frac{1}{k} \sum_{i=1}^{k} \mathbf{h}_l(t_i)
  \label{eq:meanpool}
\end{equation}
When a word appears in multiple sentences, we further average the per-sentence
vectors to produce a single representation per layer.

\subsection{Gender Geometry Construction}

For each transformer layer $l$, we build a \emph{gender subspace} from the
anchor words. For each male anchor $a^M_j$ (resp.\ female anchor $a^F_j$), we
collect 12 corpus sentences and compute its layer-$l$ representation via
Equation~\ref{eq:meanpool}. The male and female centroids are:
\begin{equation}
  \mathbf{m}_l = \frac{1}{|A^M|} \sum_{j} \mathbf{s}_l(a^M_j), \quad
  \mathbf{f}_l = \frac{1}{|A^F|} \sum_{j} \mathbf{s}_l(a^F_j)
\end{equation}
where $|A^M|=11$ and $|A^F|=10$. The unit gender direction is:
\begin{equation}
  \mathbf{g}_l = \frac{\mathbf{m}_l - \mathbf{f}_l}
                      {\|\mathbf{m}_l - \mathbf{f}_l\|}
\end{equation}
Positive projection onto $\mathbf{g}_l$ indicates proximity to the male cluster;
negative indicates proximity to the female cluster.

\subsection{Bias Estimation}

For each profession word $p$, we collect 20 corpus sentences and compute
$\mathbf{s}_l(p)$ via Equation~\ref{eq:meanpool}. We use two independent
estimators:

\paragraph{Projection bias (ProjBias).}
\begin{equation}
  \text{ProjBias}_l(p) = \mathbf{s}_l(p) \cdot \mathbf{g}_l
  \label{eq:proj}
\end{equation}
The scalar projection of the (un-normalised) profession vector onto the unit
gender direction. Preserves vector magnitude.

\paragraph{Cosine difference (CosDiff).}
\begin{equation}
  \text{CosDiff}_l(p) =
    \cos\!\left(\mathbf{s}_l(p),\,\mathbf{m}_l\right) -
    \cos\!\left(\mathbf{s}_l(p),\,\mathbf{f}_l\right)
  \label{eq:cosdiff}
\end{equation}
A purely angular measure, invariant to vector magnitude. Positive values
indicate the profession is angularly closer to the male centroid.

Both estimators are computed for every layer $l \in \{0,\ldots,12\}$ and for
all 30 professions. Their layer-wise Spearman rank correlation tests metric
agreement (\textbf{RQ3}).

\subsection{Representation Uniqueness Analysis}

To assess whether contextual embeddings genuinely encode sentence-level context,
we compute the \emph{mean pairwise cosine similarity} across all sentence
vectors. For a profession appearing in $N$ sentences with layer-$l$ vectors
$\mathbf{v}^{(1)}_l, \ldots, \mathbf{v}^{(N)}_l$:
\begin{equation}
  \overline{\cos}_l(p) = \frac{2}{N(N-1)}
    \sum_{i < j} \cos\!\left(\mathbf{v}^{(i)}_l,\,\mathbf{v}^{(j)}_l\right)
\end{equation}
A value near 1.0 indicates near-static representations; a lower value indicates
context-sensitivity. To disentangle corpus template effects from model-intrinsic
behaviour, we also ran the analysis on five hand-crafted sentences for
{\ethiopicfont ሓኪም} (doctor) spanning maximally diverse syntactic roles (subject/agent,
object/recipient, negated goal, embedded relative clause, hypothetical
conditional).

\subsection{Static Embedding Analysis}

As a complementary probe, we extract each profession word's \emph{static
embedding} directly from the token embedding weight matrix
$\mathbf{W} \in \mathbb{R}^{V \times 768}$, with no forward pass:
\begin{equation}
  \mathbf{e}(p) = \frac{1}{k}\sum_{i=1}^{k} \mathbf{W}[t_i]
\end{equation}
We build an analogous static gender direction from anchor static embeddings and
project each profession's static vector onto it. Comparing static bias scores
with contextual bias scores reveals whether gender association is encoded in the
model weights themselves or emerges during transformer forward propagation.

\subsection{Batch Stability Validation}

To verify that bias estimates are robust to sentence sampling, we partition
each profession's 499 corpus sentences into 24 non-overlapping batches of 20
sentences each (19 sentences are unused as a leftover). We compute the mean
projection score independently for each batch and report the grand mean,
standard deviation, minimum, and maximum across the 24 batches.


% ══════════════════════════════════════════════════════════════════════════════
\section{Results}
% ══════════════════════════════════════════════════════════════════════════════

\subsection{Representation Uniqueness}

Table~\ref{tab:mean_cos} reports the mean pairwise cosine similarity
$\overline{\cos}$ (averaged across all 13 layers) for every profession word,
computed from 50 sentences per profession.

\begin{table}[h]
\centering
\small
\begin{tabular}{lcc}
\toprule
\textbf{Profession} & \textbf{$\overline{\cos}$ (all layers)} & \textbf{Subword tokens} \\
\midrule
{\ethiopicfont ኣካውንታንት}  (accountant)    & 0.9809 & 5 \\
{\ethiopicfont ሰራሕተኛ ጽሬት} (cleaner)      & 0.9812 & 6 \\
{\ethiopicfont ኣርኪቴክት}   (architect)     & 0.9803 & 5 \\
{\ethiopicfont ኤለክትሪሻን}  (electrician)   & 0.9782 & 6 \\
{\ethiopicfont ዳይረክተር}   (director)      & 0.9772 & 5 \\
{\ethiopicfont ሓላዊ ጸጥታ}  (security guard) & 0.9775 & 5 \\
\multicolumn{3}{c}{\ldots (22 professions, all $\overline{\cos} > 0.94$) \ldots} \\
{\ethiopicfont ጠበቓ}      (lawyer)        & 0.9196 & 3 \\
\midrule
\textbf{Mean (all 30)}   & \textbf{0.969} & --- \\
\textbf{Min}             & \textbf{0.920} ({\ethiopicfont ጠበቓ}) & --- \\
\textbf{Max}             & \textbf{0.981} ({\ethiopicfont ሰራሕተኛ ጽሬት}) & --- \\
\bottomrule
\end{tabular}
\caption{Mean pairwise cosine similarity across sentence vectors (50 sentences,
all 13 layers). All 30 professions exceed 0.91, indicating near-static
contextual representations. {\ethiopicfont ጠበቓ} (lawyer) is the most
context-sensitive profession in the set.}
\label{tab:mean_cos}
\end{table}

All 30 professions exceed $\overline{\cos} = 0.91$. For {\ethiopicfont ሓኪም} (doctor)
specifically, $\overline{\cos} = 0.967$ across all 499 corpus sentences
($124{,}251$ pairwise comparisons at layer~12 alone; 100\% of pairs have
cosine~$> 0.99$ at that layer).

To test whether this is a corpus-template artifact, we ran the same analysis
on five hand-crafted, syntactically diverse sentences for {\ethiopicfont ሓኪም}.
The result was $\overline{\cos} = 0.947$ — a reduction of only 0.020 compared
to the template corpus — confirming that near-static representations are
intrinsic to the model's handling of Tigrigna tokens, not an artifact of
repetitive sentence structure.

\subsection{Layer-wise Projection Curve}

Figure~\ref{fig:layer_curve} shows the mean projection score across all 30
professions at each transformer layer.

\begin{figure}[h]
  \centering
  \includegraphics[width=\linewidth]{ti_projection_curve.png}
  \caption{Layer-wise mean projection score (ProjBias) averaged across 30
  Tigrigna professions on XLM-RoBERTa-base (13 layers). Positive values
  indicate male-leaning bias. The curve is near-neutral at layer~0 ($+0.20$),
  dips to $-0.62$ at layer~2, rises sharply to a single peak of $+2.99$ at
  layer~9, then collapses to $+0.04$ at layer~12. XLM-RoBERTa-large (25 layers)
  produces an analogous double-peak curve with peaks at layers~13 and~23.}
  \label{fig:layer_curve}
\end{figure}

The curve reveals a consistent three-phase structure:
\begin{enumerate}
  \item \textbf{Embedding phase (layers 0--4):} Near-neutral to slightly
    female-leaning ($-0.62$ trough at layer~2), reflecting surface token
    encoding.
  \item \textbf{Semantic accumulation phase (layers 5--9):} Monotonic rise from
    $+0.52$ to a peak of $+\mathbf{2.99}$ at \textbf{layer~9}, where attention
    heads build semantic representations and gender associations accumulate.
  \item \textbf{Output collapse phase (layers 10--12):} Sharp reduction to
    $+0.04$ at layer~12, consistent with the representation uniqueness finding
    that layer~12 produces fully collapsed, near-anisotropic vectors
    ($\overline{\cos}_{12} \approx 0.997$ for {\ethiopicfont ሓኪም}).
\end{enumerate}

\subsection{Metric Agreement — Spearman Correlation}

Table~\ref{tab:spearman} reports the Spearman rank correlation $\rho$ between
ProjBias and CosDiff across all 30 professions at each layer.

\begin{table}[h]
\centering
\small
\begin{tabular}{ccc}
\toprule
\textbf{Layer} & \textbf{Spearman $\rho$} & \textbf{$p$-value} \\
\midrule
 0 & 0.9947 & $< 10^{-28}$ \\
 1 & 0.9929 & $< 10^{-27}$ \\
 2 & 0.9933 & $< 10^{-27}$ \\
 3 & 0.9938 & $< 10^{-27}$ \\
 4 & 0.9929 & $< 10^{-27}$ \\
 5 & 0.9978 & $< 10^{-33}$ \\
 6 & 0.9969 & $< 10^{-31}$ \\
 7 & 0.9924 & $< 10^{-26}$ \\
 8 & 0.9947 & $< 10^{-28}$ \\
 9 & 0.9956 & $< 10^{-29}$ \\
10 & 0.9987 & $< 10^{-36}$ \\
11 & 1.0000 & $< 10^{-40}$ \\
12 & 1.0000 & $< 10^{-40}$ \\
\bottomrule
\end{tabular}
\caption{Spearman rank correlation between ProjBias and CosDiff across 30
professions at each layer ($n=30$). The minimum $\rho$ is $0.992$ (layer~7);
both metrics agree perfectly at layers~11--12 ($\rho=1.00$). All $p < 10^{-26}$.}
\label{tab:spearman}
\end{table}

The minimum Spearman $\rho$ across all layers is $0.992$, rising to a perfect
$\rho = 1.00$ at layers~11 and~12. This near-perfect agreement at every layer
validates that both metrics measure the same underlying bias structure and that
the choice of estimator does not affect the rank ordering of professions.

\subsection{Profession-level Bias Analysis}

Table~\ref{tab:prof_bias} summarises the mean ProjBias (averaged across all 13
layers) for key professions, alongside their static embedding bias scores.

\begin{table}[h]
\centering
\small
\begin{tabular}{lcccc}
\toprule
\textbf{Profession} & \textbf{Static proj} & \textbf{Contextual mean} & \textbf{Gain} & \textbf{Mechanism} \\
\midrule
{\ethiopicfont መምህር} (teacher)   & $+0.499$ & $+2.084$ & $+1.585$ & Embedding-level \\
{\ethiopicfont ፖሊስ}  (police)    & $+0.428$ & ---      & ---      & Embedding-level \\
{\ethiopicfont ሓረስታይ} (farmer)  & $+0.199$ & $+1.123$ & $+0.924$ & Attention-level \\
{\ethiopicfont ሓኪም}  (doctor)   & $+0.012$ & $+0.640$ & $+0.628$ & Attention-level \\
{\ethiopicfont ጠበቓ}  (lawyer)   & $-0.259$ & $-0.103$ & ---      & Consistent neutral \\
{\ethiopicfont ነርስ}   (nurse)    & $-0.352$ & $+2.748^\dagger$ & $+3.100$ & Attention-level (reversal) \\
{\ethiopicfont ነዳቓይ} (miner)    & $-0.489$ & ---       & ---      & Statically female \\
{\ethiopicfont ሳይንቲስት} (scientist) & $-0.508$ & ---    & ---      & Statically female \\
\bottomrule
\end{tabular}
\caption{Static vs.\ contextual mean projection bias for selected professions.
Gain = contextual $-$ static. $^\dagger$Nurse contextual value is the mean
ProjBias across layers 0--12; the peak at layer~9 is $+4.29$.}
\label{tab:prof_bias}
\end{table}

\subsection{Model Comparison}

Table~\ref{tab:model_compare} summarises key metrics across all five models.
M:F:N counts the number of professions classified as male-leaning, female-leaning,
and neutral (grand mean $> +0.3$, $< -0.3$, or between) across all scored layers.

\begin{table}[h]
\centering
\small
\begin{tabular}{lcccccc}
\toprule
\textbf{Model} & \textbf{Layers} & \textbf{Vocab} & \textbf{Peak layer} & \textbf{Peak proj} & \textbf{M:F:N} & \textbf{Min $\rho$} \\
\midrule
mBERT              & 13 & WP $\sim$120K & ---  & $\approx -3.21^{*}$ & n/a     & n/a \\
XLM-R-base         & 13 & SP 250K       & 9    & $+2.99$             & 26:3:1  & 0.992 \\
XLM-R-large        & 25 & SP 250K       & 23   & $+3.46$             & 25:3:2  & 0.992 \\
XLM-V-base         & 13 & SP 1M         & 11   & $+2.63$             & $25:2:1^{\dagger}$ & $0.970^{\ddagger}$ \\
mDeBERTa-v3-base   & 13 & SP            & 11   & $-1.84^{\S}$        & 2:28:0  & 0.974 \\
\bottomrule
\end{tabular}
\caption{Five-model comparison on the Tigrigna gender bias pipeline.
WP = WordPiece; SP = SentencePiece.
$^{*}$mBERT projects all 30 professions identically ($\approx -3.21$) because
every Tigrigna token maps to \texttt{[UNK]}.
$^{\dagger}$28 professions scored; 2 skipped due to tokenisation mismatch with corpus.
$^{\ddagger}$Minimum $\rho$ excluding the final-layer collapse at layer~12
($\rho = 0.72$ at layer~12 where magnitude-based and angular metrics diverge).
$^{\S}$mDeBERTa peak is female-leaning (negative); the magnitude at the
pre-collapse peak layer~11 is $-1.84$.}
\label{tab:model_compare}
\end{table}

\paragraph{mBERT — negative control.}
mBERT's WordPiece vocabulary contains no Ethiopic characters. Every one of the
30 profession words, all 21 anchor words, and every corpus sentence token maps
to the single \texttt{[UNK]} token. As a result, all profession representations
are identical at every layer, the CosDiff estimator yields $\approx 0.0002$ for
all professions (effectively zero), and ProjBias collapses to a fixed value of
$\approx -3.21$ driven entirely by the geometry of the anchor
\texttt{[UNK]} vectors. mBERT cannot meaningfully measure gender bias in
Tigrigna: vocabulary coverage is a prerequisite.

\paragraph{XLM-RoBERTa-large vs base — capacity effects.}
Both XLM-RoBERTa models produce meaningful, profession-differentiated bias
scores with Spearman $\rho \geq 0.992$ at every layer. The large model differs
from base in three structural ways:
\begin{enumerate}
  \item \textbf{Stronger embedding-level bias.} The large model's layer-0 mean
    projection is $+1.29$, versus $+0.20$ for base — a 6.5$\times$ increase in
    male bias already present before any attention layer runs.
  \item \textbf{Double-peak structure.} Base exhibits a single bias accumulation
    phase peaking at layer~9. Large produces two distinct peaks: layer~13
    ($+2.99$) and layer~23 ($+3.46$), separated by a partial suppression phase
    (layers 14--22 oscillate between $+0.91$ and $+2.74$). The extra 12
    transformer layers create a second bias amplification stage.
  \item \textbf{Final-layer behaviour.} Base collapses to $+0.04$ at layer~12;
    large retains $-0.51$ at layer~24, suggesting the large model does not
    fully suppress bias at output.
\end{enumerate}

Table~\ref{tab:prof_large} shows per-profession mean ProjBias for the large model
alongside selected base values.

\begin{table}[h]
\centering
\small
\begin{tabular}{lccl}
\toprule
\textbf{Profession} & \textbf{Base mean} & \textbf{Large mean} & \textbf{Note} \\
\midrule
{\ethiopicfont ዳኛ}        (judge)   & ---      & $+3.67$ & Strongest male in large \\
{\ethiopicfont መምህር}      (teacher) & $+2.08$  & $+3.16$ & Amplified by capacity \\
{\ethiopicfont ሓላፊ}       (manager) & ---      & $+2.54$ & Strong male \\
{\ethiopicfont ፖሊስ}       (police)  & ---      & $+2.96$ & Strong male \\
{\ethiopicfont ሓኪም}       (doctor)  & $+0.64$  & $+1.17$ & Consistent male, stronger in large \\
{\ethiopicfont ሓረስታይ}     (farmer)  & $+1.12$  & $+1.85$ & Consistent male \\
{\ethiopicfont ነርስ}        (nurse)   & $+2.75$  & $+0.47$ & Male in both; weaker in large \\
{\ethiopicfont ጠበቓ}       (lawyer)  & $-0.10$  & $-0.81$ & Neutral in base → female in large \\
{\ethiopicfont ሳይንቲስት}    (scientist)& ---     & $-0.86$ & Female in large \\
{\ethiopicfont መካኒክ}      (mechanic) & ---     & $-0.91$ & Female in large \\
\bottomrule
\end{tabular}
\caption{Mean ProjBias (averaged across all layers) for XLM-RoBERTa-base and
XLM-RoBERTa-large. Base mean omitted where not previously reported.
{\ethiopicfont ጠበቓ} (lawyer) is neutral in base but female-leaning in large,
illustrating how increased model capacity can shift borderline professions.}
\label{tab:prof_large}
\end{table}

The most striking individual shift is {\ethiopicfont ጠበቓ} (lawyer): base assigns it a
grand mean of $-0.10$ (neutral, all 24 batches neutral), while large assigns
$-0.81$ (female-leaning). The large model's deeper representation space
apparently resolves the ambiguity that base leaves unresolved.

\paragraph{XLM-V-base — consistent with XLM-R family.}
XLM-V-base produces bias patterns consistent with the XLM-R family for its
28 scored professions: 25 male-leaning, 2 female-leaning, 1 neutral, with
a peak mean projection of $+2.63$ at layer~11. The layer-wise curve resembles
XLM-R-base in shape but peaks one layer later (layer~11 vs layer~9) and
exhibits a more severe final-layer collapse ($-3.27$ vs $+0.04$). Spearman
$\rho \geq 0.970$ for layers 0--11, confirming consistent metric agreement
before the collapse. The 4$\times$ vocabulary increase over XLM-R does not
alter the direction of measured bias.

\paragraph{mDeBERTa-v3-base — complete bias direction reversal.}
mDeBERTa-v3-base produces a striking reversal: 28 of 30 professions are
female-leaning (mean projection $< -0.3$), with only {\ethiopicfont ፕሮፌሰር}
(professor, grand mean $+0.27$) and {\ethiopicfont ሓረስታይ} (farmer, $+0.13$)
remaining male-leaning. The layer-wise mean projection is consistently negative
from layer~5 onwards ($-0.97$ at layer~5, $-1.84$ at layer~11), with a sharp
final-layer collapse to $-7.24$ at layer~12 — far more extreme than any other
model's final layer. Spearman $\rho \geq 0.974$ at every layer confirms the
reversal is internally consistent across both estimators. This finding is
discussed in detail in Section~\ref{sec:deberta_reversal}.

\subsection{Batch Stability Validation}

Table~\ref{tab:batch} reports the grand mean, standard deviation, and range
of the per-batch mean projection score across 24 batches of 20 sentences each.

\begin{table}[h]
\centering
\small
\begin{tabular}{lccccc}
\toprule
\textbf{Profession} & \textbf{Grand mean} & \textbf{Std dev} & \textbf{Min} & \textbf{Max} & \textbf{M/F/N batches} \\
\midrule
{\ethiopicfont መምህር} (teacher)  & $+2.084$ & $0.045$ & $+2.009$ & $+2.181$ & 24 / 0 / 0 \\
{\ethiopicfont ሓረስታይ} (farmer) & $+1.123$ & $0.067$ & $+0.998$ & $+1.248$ & 24 / 0 / 0 \\
{\ethiopicfont ሓኪም}  (doctor)  & $+0.640$ & $0.092$ & $+0.529$ & $+0.850$ & 24 / 0 / 0 \\
{\ethiopicfont ጠበቓ}  (lawyer)  & $-0.103$ & $0.089$ & $-0.230$ & $+0.091$ &  0 / 0 / 24 \\
\bottomrule
\end{tabular}
\caption{Batch stability across 24 non-overlapping batches of 20 sentences each.
M/F/N counts show how many batches yielded a MALE / FEMALE / NEUTRAL verdict.
All 24 batches for teacher, farmer, and doctor are MALE-LEANING regardless of
the local gender-sentence ratio.}
\label{tab:batch}
\end{table}

For {\ethiopicfont መምህር} (teacher), the standard deviation is only $0.045$ — the
batch estimates vary by less than 2\% of the grand mean. The floor for teacher
($+2.009$) is higher than the ceiling for doctor ($+0.850$), placing these two
professions in entirely non-overlapping bias ranges. For
{\ethiopicfont ጠበቓ} (lawyer), no batch crosses the $\pm 0.3$ threshold in either
direction (range $[-0.230, +0.091]$).


% ══════════════════════════════════════════════════════════════════════════════
\section{Observations and Discussion}
% ══════════════════════════════════════════════════════════════════════════════

\subsection{Near-static Representations in a Low-Resource Language}

Our most foundational finding is that XLM-RoBERTa-base produces effectively
static representations for all 30 Tigrigna profession tokens. The mean pairwise
cosine similarity exceeds $0.91$ for every profession (Table~\ref{tab:mean_cos}),
indicating that the model assigns nearly identical hidden-state vectors to a
given token regardless of surrounding discourse context.

This behaviour is not primarily a corpus artifact. The custom-sentence
experiment for {\ethiopicfont ሓኪም}, using five structurally and semantically maximally
diverse sentences, produced $\overline{\cos} = 0.947$ compared to $0.967$ from
the template corpus — a reduction of only $0.020$. Near-static representations
are therefore intrinsic to the model's treatment of Tigrigna tokens.

We attribute this to the low-resource nature of Tigrigna in XLM-RoBERTa's
pretraining corpus. When a language is sparsely represented, the model's
attention heads lack the statistical signal to learn context-sensitive
representations, causing the transformer to degenerate toward a near-lookup-table
behaviour where representations are driven primarily by token identity rather
than surrounding context.

\subsection{Two Distinct Mechanisms of Gender Bias}

Comparing static embedding bias scores with contextual bias scores reveals two
qualitatively distinct mechanisms:

\paragraph{Type 1: Embedding-level bias.}
For {\ethiopicfont መምህር} (teacher, static $+0.499$, contextual $+2.084$), the male
association is already encoded in the raw token embedding before any attention
is applied. The transformer layers then amplify this signal more than fourfold.
The bias originates in the pretraining data distribution.

\paragraph{Type 2: Attention-level bias.}
For {\ethiopicfont ሓኪም} (doctor, static $+0.012$) and {\ethiopicfont ሓረስታይ} (farmer, static
$+0.199$), the raw token embeddings are effectively gender-neutral. Yet the
transformer layers impose strong male bias ($+0.640$ and $+1.123$ respectively)
during forward propagation. The attention mechanism introduces gender
associations absent from the input representation, presumably by activating
co-occurrence patterns learned from higher-resource languages in pretraining.

\paragraph{The nurse reversal.}
The most dramatic instance is {\ethiopicfont ነርስ} (nurse), which has a statically
\emph{female}-leaning embedding (static $= -0.352$) but emerges from the
transformer with a strongly \emph{male}-leaning contextual representation,
peaking at $+4.29$ at layer~9 — the highest single-profession value in our
dataset. The transformer not only overrides but completely reverses the static
signal, suggesting that the multilingual pretraining signal for \textit{nurse}
in higher-resource languages (where the model has richer data) dominates the
Tigrigna token's own embedding.

\subsection{Layer-wise Bias Structure}

The layer-wise projection curve (Figure~\ref{fig:layer_curve}) exhibits a
consistent three-phase pattern: near-neutral early layers reflecting
morphological processing; a monotonic accumulation of male bias through layers
5--9 peaking at $+2.99$; and a collapse to near-zero at layer~12. Layer~9 is
the critical layer for gender bias in XLM-RoBERTa on Tigrigna, with individual
profession scores ranging from $+5.90$ ({\ethiopicfont መምህር}) to negative values for
female-leaning professions, the ordering preserved near-perfectly across all
layers (Spearman $\rho \geq 0.992$).

\subsection{Lawyer as a Contrast Case}

{\ethiopicfont ጠበቓ} (lawyer) is the single profession that remains consistently neutral
across all analyses. It has the lowest mean pairwise cosine similarity of all
30 professions ($\overline{\cos} = 0.920$, most contextually sensitive), the
most female-adjacent static embedding among neutral-labelled words (static
$= -0.259$), yet the transformer does not amplify this signal. All 24 sampling
batches return a ROUGHLY NEUTRAL verdict (grand mean $= -0.103$, range
$[-0.230, +0.091]$). The inverse relationship between $\overline{\cos}$ and
bias magnitude — highest context-sensitivity paired with lowest contextual bias
— is consistent with the hypothesis that contextual diversity suppresses bias
amplification.

\subsection{Vocabulary Coverage as a Prerequisite for Bias Measurement}

The mBERT experiment provides a clean negative control: a model that scores well
on standard multilingual benchmarks but is entirely blind to Tigrigna text.
Because Ethiopic script is absent from mBERT's WordPiece vocabulary, the model
cannot distinguish \ti{ሓኪም} (doctor) from \ti{ጠበቓ} (lawyer) — both are
\texttt{[UNK]}. Any bias pipeline applied to such a model would produce
numerically non-trivial but entirely meaningless results, reflecting only the
geometry of the \texttt{[UNK]} anchor embeddings.

This finding has a practical implication: before applying a multilingual bias
probe to a new language, researchers should first verify that the model's
tokeniser produces non-\texttt{[UNK]} subwords for the target vocabulary.
For low-resource and non-Latin-script languages, this is not guaranteed even
for models marketed as ``multilingual''. SentencePiece models with large
vocabularies ($\geq 128{,}000$ pieces) and character-level fallback — such as
both XLM-RoBERTa variants — are substantially more likely to provide meaningful
token-level representations.

\subsection{Capacity Effects: Base vs Large Architecture}

Scaling from XLM-RoBERTa-base (12 layers, 768-dim) to large (24 layers, 1024-dim)
does not change the \emph{direction} of gender bias — both models assign male
bias to the same broad set of professions — but it does change the
\emph{magnitude}, \emph{structure}, and \emph{resolution} of bias representation.

The double-peak structure in the large model (peaks at layers 13 and 23)
suggests that increased depth creates two semi-independent bias accumulation
circuits: an early-to-mid-network phase analogous to the single phase in base,
and a late-network phase with no counterpart in the shallower model. Whether
these two phases correspond to distinct linguistic computations (e.g.,
semantic vs.\ pragmatic) is an open question requiring attention head analysis.

The most diagnostically useful capacity effect is the resolution of borderline
cases. {\ethiopicfont ጠበቓ} (lawyer) is the clearest example: base cannot determine
its bias direction (24 batches, all neutral), while large assigns it a
consistent $-0.81$ (female-leaning). Increased representational capacity
apparently provides enough degrees of freedom to resolve gender associations
that base leaves ambiguous. This suggests that larger models may be
\emph{more biased in aggregate} — not because they introduce new associations,
but because they resolve associations that smaller models encode only weakly.

\subsection{Disentangled Attention Reverses Bias Direction}
\label{sec:deberta_reversal}

The most surprising finding of our five-model comparison is that
mDeBERTa-v3-base assigns \emph{female}-leaning bias to 28 of 30 Tigrigna
professions — the complete opposite of all three XLM-R family models, which
assign \emph{male}-leaning bias to 25--26 of the same professions. The only
two professions that remain male-leaning under mDeBERTa are
{\ethiopicfont ፕሮፌሰር} (professor, grand mean $+0.27$) and {\ethiopicfont ሓረስታይ}
(farmer, $+0.13$), both occupations with strong male associations across many
languages and cultures.

The mDeBERTa layer-wise mean projection curve is consistently negative from
layer~5 onwards, deepening monotonically from $-0.97$ at layer~5 to $-1.84$
at layer~11, with a sharp final-layer collapse to $-7.24$ at layer~12. This
collapse is far more extreme than the corresponding behaviour in XLM-R models
($+0.04$ for base, $-3.27$ for XLM-V-base). The Spearman rank correlation
between ProjBias and CosDiff remains $\rho \geq 0.974$ at every layer
($p < 10^{-19}$), confirming that the reversal is internally consistent and
not a magnitude artifact: both estimators agree that the same set of professions
cluster closer to the female anchor centroid.

We attribute this to the Disentangled Attention mechanism~\cite{he2021deberta}.
Standard self-attention (BERT/RoBERTa) builds query, key, and value
representations by summing content and positional embeddings \emph{before}
the attention computation. Disentangled Attention encodes content and position
\emph{separately}, producing different cross-token interaction patterns during
pretraining. For Tigrigna tokens — which are processed as subwords with limited
pretraining context — this appears to produce a gender subspace in which
profession vectors cluster near the female centroid rather than the male centroid.

A key question is whether this reflects genuinely learned female associations
for these professions, or an inverted geometric relationship between anchor
clusters. The CosDiff metric (purely angular, centroid-relative) shows that
professions are angularly closer to the female anchor cluster than to the male
anchor cluster — which is consistent with learned association rather than a
simple magnitude inversion. However, a definitive answer would require targeted
intervention experiments (e.g., probing with gender-marked Tigrigna sentences
or attention pattern analysis) beyond the scope of the present work.

This finding carries a methodological warning: researchers applying a single
bias probe design across multiple model \emph{families} may obtain
directionally opposite results without realising it. A profession classified
as male-biased under an XLM-R model may appear female-biased under mDeBERTa
on the same corpus. Reporting bias direction alongside bias magnitude, and
validating with at least two architectural families, is therefore important
for reproducibility.

\subsection{XLM-V Tokenisation Coverage and Corpus Alignment}

XLM-V-base's 1{,}000{,}000-piece SentencePiece vocabulary achieves broad
Ethiopic character coverage for 28 of 30 target professions. The two excluded
professions — {\ethiopicfont ጠበቓ} (lawyer) and {\ethiopicfont ነዳቓይ} (miner) — tokenise
differently under the 1M vocabulary compared to the 250K XLM-RoBERTa vocabulary
used to build our corpus. Because our context-retrieval step locates profession
words via string matching against corpus sentences, a tokeniser that segments a
word into different subword pieces may fail to locate any matching context
sentences for that word.

This reveals a \emph{vocabulary-corpus alignment} requirement: when applying a
bias pipeline to a new model, one must verify that the tokeniser's segmentation
of target vocabulary items is compatible with the corpus construction process.
The 4$\times$ vocabulary increase in XLM-V improves Ethiopic coverage in
principle but creates a mismatch with corpora built for a different tokenisation
scheme. Future work should rebuild corpora using each model's own tokeniser,
or use character-level (rather than subword-level) matching for context
retrieval to ensure portability across vocabularies.

For the 28 professions successfully scored, XLM-V-base's bias profile is
consistent with the XLM-R family: 25 male-leaning, 2 female-leaning, 1 neutral,
with a peak mean projection of $+2.63$ at layer~11. The 4$\times$ vocabulary
increase does not alter the direction of measured bias, confirming that this
aspect of the XLM-R architecture family is robust to vocabulary scaling.

\subsection{Implications for Low-Resource Language Bias Research}

\begin{enumerate}
  \item \textbf{Verify tokeniser coverage before probing.} Our mBERT
    experiment shows that a multilingual model can produce numerically
    non-trivial but entirely uninformative bias scores when the target language
    is absent from its vocabulary. Tokeniser coverage should be validated as a
    mandatory first step for any multilingual bias study on low-resource or
    non-Latin-script languages.

  \item \textbf{Contextual bias measurement may be unreliable for low-resource
    languages.} When $\overline{\cos} > 0.95$, averaging over multiple sentence
    vectors adds negligible information; the bias score is effectively being read
    from a single vector, equivalent to static embedding probing.

  \item \textbf{Batch stability validates measurement reliability.} Standard
    deviations as low as $0.045$ across 24 batches demonstrate that 20
    sentences already produce fully converged bias estimates for near-static
    representations, providing a principled basis for small context windows in
    low-resource settings.

  \item \textbf{The two-mechanism framework enables targeted debiasing.}
    Type~1 bias (embedding-level) may be addressable via gender-neutral
    retrofitting of the static weight matrix. Type~2 bias (attention-level)
    requires intervention in the transformer computation — for example,
    projection-based nullspace methods applied to the peak bias layer (layer~9
    in XLM-R-base, layer~23 in large).

  \item \textbf{Larger models amplify and resolve bias.} XLM-RoBERTa-large
    produces stronger overall male bias ($+3.46$ peak vs $+2.99$) and resolves
    ambiguous professions (lawyer: neutral in base, female-leaning in large).
    Scaling multilingual models may increase bias precision without necessarily
    changing bias direction.

  \item \textbf{Perfect metric agreement simplifies reporting.} Spearman
    $\rho \geq 0.970$ between ProjBias and CosDiff at every layer, in all
    four meaningful models, means both metrics are interchangeable for ranking
    professions by bias magnitude regardless of model scale or architecture
    family.

  \item \textbf{Attention architecture determines bias direction.} XLM-R-family
    models (standard self-attention) consistently assign male-leaning bias to
    Tigrigna professions, while mDeBERTa (Disentangled Attention) assigns
    female-leaning bias to the same professions. Bias direction is not a
    universal property of the training data alone — it depends on the
    architectural mechanism through which the model processes cross-lingual
    representations. Multi-architecture replication is therefore essential for
    robust multilingual bias conclusions.
\end{enumerate}


% ── References ────────────────────────────────────────────────────────────────
\begin{thebibliography}{99}

\bibitem{conneau2020xlm}
Conneau, A., Khandelwal, K., Goyal, N., Chaudhary, V., Wenzek, G.,
Guzm\'{a}n, F., Grave, \'{E}., Ott, M., Zettlemoyer, L., \& Stoyanov, V.
(2020).
Unsupervised cross-lingual representation learning at scale.
In \textit{Proceedings of ACL 2020}, pp.\ 8440--8451.

\bibitem{he2021deberta}
He, P., Liu, X., Gao, J., \& Chen, W. (2021).
DeBERTa: Decoding-enhanced BERT with disentangled attention.
In \textit{Proceedings of ICLR 2021}.

\bibitem{liang2023xlmv}
Liang, D., Gonen, H., Mao, Y., Hou, R., Goyal, N., Ghazvininejad, M.,
Zettlemoyer, L., \& Levy, O. (2023).
XLM-V: Overcoming the vocabulary bottleneck in multilingual masked language models.
In \textit{Proceedings of EMNLP 2023}, pp.\ 13142--13152.

\bibitem{bolukbasi2016man}
Bolukbasi, T., Chang, K.-W., Zou, J., Saligrama, V., \& Kalai, A. (2016).
Man is to computer programmer as woman is to homemaker?
Debiasing word embeddings.
In \textit{Advances in NeurIPS}, vol.\ 29.

\bibitem{zhao2019gender}
Zhao, J., Wang, T., Yatskar, M., Cotterell, R., Ordonez, V., \& Chang, K.-W.
(2019).
Gender bias in contextualized word embeddings.
In \textit{Proceedings of NAACL 2019}, pp.\ 629--634.

\bibitem{guo2021detecting}
Guo, W., \& Caliskan, A. (2021).
Detecting emergent intersectional biases: Contextualized word embeddings
contain a distribution of human-like biases.
In \textit{Proceedings of AIES 2021}, pp.\ 122--133.

\bibitem{rogers2020primer}
Rogers, A., Kovaleva, O., \& Rumshisky, A. (2020).
A primer in BERTology: What we know about how BERT works.
\textit{Transactions of ACL}, 8, 842--866.

\end{thebibliography}

\end{document}
